% Partie : théorie de la démonstration
% Chapitre : déduction naturelle
% Section : introduction

\documentclass[../../../include/open-logic-section]{subfiles}

\begin{document}

\olfileid[fr]{pt}{nat}{int}

\olsection{Introduction}

La déduction naturelle est une famille de systèmes de preuve
dans lesquels chaque connecteur logique ou quantificateur possède
deux types d'inférence : une règle qui l'introduit et une règle
qui l'élimine. Par exemple, \Intro{\lif} conclut à une formule
de la forme~$!A \lif !B$, tandis que \Elim{\lif} a une formule
de la forme~$!A \lif !B$ parmi ses prémisses.

Les deux premières versions de la déduction naturelle ont été
introduites dans les années 1930 par Gerhard Gentzen et
Stanisław Jaśkowski. Le système de Gentzen est celui qu'étudie
le plus souvent la théorie de la démonstration ; celui de Jaśkowski
est à l'origine des systèmes les plus employés dans l'enseignement.
Dans les deux cas, on peut partir d'\emph{hypothèses} :
des formules qui n'ont pas à être démontrées, mais qui peuvent
servir à en démontrer d'autres. Une preuve entière peut dépendre
de telles hypothèses. Comme celles-ci ne sont pas démontrées,
la preuve n'établit pas sa conclusion, ou \emph{formule finale},
indépendamment d'elles : elle établit que cette conclusion
en découle.

Certaines règles permettent de conclure sans dépendre de certaines
hypothèses : elles les \emph{déchargent}. Ainsi, \Intro{\lif}
prend une preuve de la conclusion~$!B$ sous l'hypothèse~$!A$
et conclut à $!A \lif !B$. La preuve se terminant par $!A \lif !B$
ne dépend plus de l'hypothèse $!A$, qui a été déchargée.
Dans les systèmes de Gentzen, les preuves sont des arbres de formules,
les hypothèses figurent en haut, et des étiquettes portées par
les règles telles que~\Intro{\lif} désignent les hypothèses
déchargées. Dans les systèmes de Jaśkowski, les parties d'une
preuve dépendant d'une hypothèse sont délimitées, par exemple
par une indentation derrière un trait vertical ou par un cadre.
L'hypothèse~$!A$ déchargée et le conséquent $!B$ de l'implication
figurent à la première et à la dernière ligne de ce cadre ;
la conclusion~$!A \lif !B$ de \Intro{\lif} se trouve à l'extérieur.

Ces systèmes sont dits naturels parce que leurs règles reflètent
les formes usuelles du raisonnement mathématique.
Pour établir une implication, on suppose son antécédent et
l'on s'en sert pour démontrer son conséquent : c'est
\Intro{\lif}. Si l'on connaît une implication $!A \lif !B$
et son antécédent~$!A$, on conclut à~$!B$ : c'est \Elim{\lif}.
Une preuve par cas permet de tirer une conclusion commune
d'une disjonction~$!A \lor !B$ : c'est \Elim{\lor}.
Pour établir un énoncé général~$\lforall[x][!A(x)]$,
on démontre $!A(c)$ pour un paramètre arbitraire~$c$,
respectant la condition de constante propre :
c'est \Intro{\lforall}. Et ainsi de suite.

Voici, par exemple, une preuve simple de~$!A \lif (!A \lor !B)$
en déduction naturelle à la Gentzen :
\[
\AxiomC{$\Discharge{!A}{x}$}
\RightLabel{\Intro{\lor}}
\UnaryInfC{$!A \lor !B$}
\DischargeRule{\Intro{\lif}}{x}
\UnaryInfC{$!A \lif (!A \lor !B)$}
\DisplayProof
\]
Elle part de l'hypothèse~$!A$, applique \Intro{\lor} pour
obtenir~$!A \lor !B$, puis \Intro{\lif} pour conclure
à~$!A \lif (!A \lor !B)$. L'inférence \Intro{\lif}
décharge l'hypothèse~$!A$, comme l'indique l'étiquette~$x$.

La théorie de la démonstration emploie surtout les systèmes
originaux de Gentzen, fondés sur des arbres de formules.
Une variante présente cependant aussi un intérêt.
Une preuve de~$!A$ sous les hypothèses~$\Gamma$ montre que
$!A$ découle de~$\Gamma$, c'est-à-dire que
$\Gamma \Entails !A$. On peut donc lire les règles de déduction
naturelle comme des règles portant sur ces conséquences logiques.
Ainsi, \Intro{\lif} exprime que, si
$\Gamma \cup \{!A\} \Entails !B$, alors
$\Gamma \Entails !A \lif !B$.

Il est naturel de formuler directement les règles de façon
qu'elles portent à la fois sur les hypothèses et sur la conclusion
qui en découle, dans leurs prémisses et leur conclusion :
elles opèrent alors sur des séquents $\Gamma \Sequent !A$.
La preuve simple ci-dessus s'écrit ainsi dans un tel système :
\[
\Axiom$x:!A \fCenter !A$
\RightLabel{\Intro{\lor}}
\UnaryInf$x:!A \fCenter !A \lor !B$
\DischargeRule{\Intro{\lif}}{x}
\UnaryInf$\fCenter !A \lif (!A \lor !B)$
\DisplayProof
\]
Les règles sont les mêmes : elles agissent sur la formule
à droite de~$\Sequent$, tandis que les formules à gauche
énumèrent les hypothèses dont chaque étape dépend.

L'ensemble des règles d'introduction et d'élimination ne suffit
pas à lui seul à la complétude pour la logique classique.
Il faut ajouter deux règles liées à~$\lfalse$.
La première est la \emph{règle d'absurdité intuitionniste}~\FalseInt,
ou \emph{explosion} : de $\lfalse$, on peut déduire n'importe
quelle formule. La seconde est le principe classique de
preuve indirecte~\FalseCl, ou \emph{règle d'absurdité classique} :
si l'on démontre~$\lfalse$ à partir de~$\lnot !A$,
on peut conclure à~$!A$ en déchargeant cette hypothèse.
Sans \FalseCl, on obtient la logique intuitionniste.
Sans ces deux règles, on obtient la \emph{logique minimale}.

Nous étudierons la déduction naturelle à la Gentzen pour
les logiques classique et intuitionniste, dans les systèmes
\Log{N1c} et \Log{N1i}. Les systèmes \Log{N2c} et \Log{N2i}
sont leurs versions correspondantes sur les séquents
$\Gamma \Sequent !A$, plutôt que sur les seules formules~$!A$.










\end{document}
